%
Adaptive bitrate (ABR) algorithms play a vital role in video streaming by adjusting video quality under varying and unpredictable network conditions. They dynamically adjust quality to maintain stable playback in various environments, such as fluctuating network conditions. However, evaluating their performance and choosing a suitable algorithm is challenging due to variations in test environments and network traces. This thesis investigates the performance and behavior of six ABR algorithms; BOLA, Linear BBA, TTP, Gelato, Pensieve, and MPC using the Puffer live video streaming platform. To maintain a fair environment for all ABR algorithms, we stream pre-recorded videos in an emulated environment using Mahimahi and CellReplay, two network emulators designed for evaluation of internet systems. We run various traces with different conditions, bandwidths, and signal strengths to reveal clear trade-offs for ABR algorithms. We collect QoE (quality of experience) and network performance-related metrics for each ABR algorithm. A key contribution of this thesis is a cross-emulator comparison. CellReplay is a recent, award-winning tool that replaced the state-of-the-art in cellular network record-and-replay. This thesis investigates how the performance and ranking of ABR algorithms change when CellReplay is used. To address this, we run the same traces in both emulators and analyze their impact on the performance of each ABR algorithm. Our analysis shows that ABR rankings are sensitive to both emulator choices and network conditions. We also compare our results with previous work to evaluate in what way our work matches to theirs. This work provides a practical and extensible framework for ABR evaluation, contributing new insights into the stability of ABR performance across testbeds and trace types.


%%% Local Variables:
%%% mode: latex
%%% TeX-master: "../thesis"
%%% End:
